\newpage
\section{Misión 1: Compartir es vivir}\label{sec:mission1}

\begin{wrapfigure}{L}{0.23\linewidth}
\vspace{-0.5cm}\includegraphics[width=\linewidth]{abacus}
\end{wrapfigure}
%
Hace unos 75 años, una expedición secreta fue enviada al espacio con el objetivo de colonizar otro planeta. La expedición llegó sana y salva a su destino, pero el aterrizaje destrozó sus instrumentos de comunicación y los dejó aislados. Tú eres un colono de tercera generación, futuro ingeniero, y esta es la primera misión que te asigna la asamblea:
compartir la nueva y reluciente \mbox{calculadora}. \textbf{Por favor, lee todo con \mbox{atención}}.\\

\noindent
\begin{remark}
\textbf{Importante}: Se anima a cada estudiante a publicar un \textbf{informe formal} (en PDF) en el ``Foro de intercambio de la Misión 1'' del Campus Virtual antes de la siguiente sesión de 2h. Después de eso, contamos contigo para publicar una \textbf{carta de respuesta formal con comentarios} (también en PDF) sobre al menos otra solución, antes de la siguiente sesión de 1h.
\end{remark}

\subsection*{Detalles del informe}

Ahora mismo, la colonia consta de 3 bases separadas (3 edificios). Estas bases tienen 3, 200 y 5000 dispositivos electrónicos, respectivamente. La base con 3 dispositivos se ha construido específicamente para alojar una calculadora nueva, potente y con gran consumo de energía (es uno de los 3 dispositivos). Todo el mundo en la colonia quiere tener acceso a la calculadora, pero ninguno de los dispositivos tiene capacidades de red \textit{todavía}. Tu objetivo es proponer un diseño de red convincente que conecte y satisfaga a todo el mundo.\\

\noindent
Prepara un breve informe con tu diseño, proponiendo soluciones claras y completas a los siguientes puntos:

\begin{enumerate}[itemsep=0.1cm]
\item ¿Cómo se cruza la \textbf{barrera analógico-digital} para permitir la comunicación digital? Sé concreto en las tecnologías que propongas.
\item ¿Cómo se conectan físicamente los dispositivos? Aporta un diagrama con la \textbf{topología} exacta.
\item ¿Qué \textbf{sistema de direccionamiento} se usará para identificar de forma única cada dispositivo? Da dos ejemplos de dirección, uno de ellos el de la nueva calculadora.
\item ¿Qué \textbf{formato de paquete} emplearán los dispositivos? Aporta un diagrama de formato concreto y los bytes de un paquete de ejemplo completo.
\end{enumerate}


\subsection*{Detalles de los comentarios}

Después de publicar tu solución en el foro, busca otra y \textbf{reclámala} respondiendo con un breve mensaje ``Yo daré mi opinión'' (una carta de comentarios por solución).\\

\noindent
Después prepara una carta formal en PDF y publícala en el mismo hilo cuando esté lista (puedes encontrar un ejemplo de carta de comentarios en el Campus Virtual). En esta carta, incluye tus comentarios (tanto positivos como negativos), explicando todo aquello que consideres relevante o que creas que puede ayudar a esa solución.  Asegúrate de tratar los siguientes puntos:

\begin{enumerate}[itemsep=0.1cm]
\item ¿Ves algún fallo importante o algún punto incompleto que impediría que la propuesta se hiciera realidad?
\item Si volvieras a hacer tu informe, ¿adoptarías alguna de las propuestas de esta solución?
\item ¿Has podido interpretar el contenido del mensaje incluido en el informe? Incluye los bytes de una posible respuesta, según tu comprensión de ese informe.
\end{enumerate}

